\chapter{Cài đặt hệ thống}

\section{Cài đặt thiết bị nhúng ESP32}

\subsection{Môi trường phát triển}

Firmware cho ESP32-CAM được phát triển trên Arduino IDE với ESP32 Board Package
của Espressif. Các thư viện chính được sử dụng bao gồm: \texttt{esp\_camera.h}
cho điều khiển camera OV5640, \texttt{driver/i2s.h} cho giao tiếp I2S với
microphone INMP441, và \texttt{WiFi.h} cho kết nối mạng. Firmware được flash
qua UART với baudrate 115200.

\subsection{Kết nối WiFi và khởi tạo}

Khi khởi động, ESP32 thực hiện tuần tự các bước: kết nối WiFi đến access point
được cấu hình sẵn, khởi tạo camera OV5640 với các tham số mặc định (QVGA,
JPEG quality 12), cấu hình I2S cho microphone INMP441, và cuối cùng thiết lập
kết nối TCP đến server trung gian. Nếu mất kết nối WiFi trong quá trình hoạt
động, ESP32 tự động thử kết nối lại với khoảng cách 5 giây giữa các lần thử.

\subsection{Cấu hình I2S}

\begin{table}[htbp]
\centering
\caption{Cấu hình I2S cho microphone INMP441}
\label{tab:i2s-config}
\begin{tabular}{>{\raggedright\arraybackslash}p{0.35\linewidth} >{\raggedright\arraybackslash}p{0.55\linewidth}}
\toprule
\rowcolor{headerblue} \textcolor{white}{\textbf{Tham số}} & \textcolor{white}{\textbf{Giá trị}} \\
\midrule
\rowcolor{rowgray} Chế độ & I2S\_MODE\_MASTER | I2S\_MODE\_RX \\
Tần số lấy mẫu & 16.000 Hz \\
\rowcolor{rowgray} Bit depth đầu vào & 32-bit (I2S native) \\
Bit depth đầu ra & 16-bit PCM (sau chuyển đổi) \\
\rowcolor{rowgray} Kênh & Mono (chỉ kênh trái) \\
Số DMA buffer & 16 (tăng từ 8 để ổn định) \\
\rowcolor{rowgray} Kích thước buffer & 1024 mẫu (64 ms mỗi buffer) \\
Throughput & 32 KB/s ($\approx$ 256 kbps) \\
\bottomrule
\end{tabular}
\end{table}

Số lượng DMA buffer được tăng từ 8 lên 16 sau quá trình thử nghiệm, nhằm
giảm hiện tượng buffer overflow khi Core 1 bận xử lý camera capture. Với 16
buffer, hệ thống có thể chịu được độ trễ tối đa 16 $\times$ 64 ms = 1024 ms
trước khi mất dữ liệu âm thanh.

\subsection{Chuẩn hóa tín hiệu}

Microphone INMP441 có độ phân giải 18-bit, đóng gói trong word 32-bit với dữ liệu
nằm ở các bit cao. Quá trình chuẩn hóa bắt đầu bằng việc đọc mẫu 32-bit raw
từ I2S DMA buffer, sau đó dịch phải 14 bit
(\texttt{val = raw\_samples[i] >> 14}) để đưa 18 bit dữ liệu hữu ích về
phạm vi 16-bit. Giá trị sau đó được cắt ngưỡng vào $[-32768, 32767]$ để tránh
tràn khi có tín hiệu mạnh bất thường, và lưu dưới dạng \texttt{int16\_t}
để truyền qua TCP.

Việc dịch phải 14 bit (thay vì 16 bit) là lựa chọn có chủ đích: giữ lại
toàn bộ 18 bit dữ liệu hữu ích thay vì mất 2 bit thấp, tận dụng tối đa
độ phân giải của microphone. Trong thực tế, 2 bit thấp nhất thường chứa
nhiễu lượng tử nên việc mất chúng qua bước cắt ngưỡng không ảnh hưởng
đáng kể đến chất lượng âm thanh.

\subsection{Pipeline streaming video}

Camera OV5640 được cấu hình chụp JPEG với double buffering (2 frame buffer)
để tối ưu throughput --- trong khi một buffer đang được gửi qua TCP, camera
có thể ghi vào buffer còn lại. XCLK mặc định 12 MHz, chất lượng JPEG có
thể điều chỉnh từ 10 (chất lượng cao nhất, file lớn) đến 63 (chất lượng
thấp nhất, file nhỏ).

Mỗi frame được truyền qua TCP port 3001 với giao thức tự định nghĩa:
\texttt{[Length (4 bytes LE)][JPEG Data (variable)]}. Header 4 byte chứa
độ dài frame theo little-endian, cho phép server biết chính xác kích thước
dữ liệu cần đọc. Buffer tĩnh 250 KB được cấp phát một lần trên PSRAM
(hoặc heap nếu PSRAM không khả dụng) để tránh phân mảnh bộ nhớ do
\texttt{malloc/free} lặp lại.

\begin{figure}[htbp]
\centering
\fbox{\parbox{0.85\linewidth}{\centering\vspace{2cm}
\textit{Thêm hình: Sơ đồ cấu trúc TCP frame --- hiển thị header 4 bytes
(Length, little-endian) theo sau bởi JPEG data có kích thước variable.
Chú thích ví dụ: Length = 0x00 0x8C 0xA0 0x00 = 36.000 bytes.}
\vspace{2cm}}}
\caption{Cấu trúc TCP frame cho video streaming}
\label{fig:tcp-frame}
\end{figure}

\subsection{Giao thức điều khiển camera}

Trình duyệt gửi lệnh điều khiển qua WebSocket đến server trung gian,
server chuyển tiếp qua TCP đến ESP32 với format:
\texttt{[0xA5][ID (1 byte)][Value (1 byte)]}. Byte magic \texttt{0xA5}
giúp ESP32 phân biệt lệnh điều khiển với dữ liệu thông thường trên
cùng kết nối TCP.

\begin{table}[htbp]
\centering
\caption{Các lệnh điều khiển camera chính}
\label{tab:camera-commands}
\begin{tabular}{>{\centering\arraybackslash}p{0.08\linewidth} >{\raggedright\arraybackslash}p{0.3\linewidth} >{\raggedright\arraybackslash}p{0.45\linewidth}}
\toprule
\rowcolor{headerblue} \textcolor{white}{\textbf{ID}} & \textcolor{white}{\textbf{Tham số}} & \textcolor{white}{\textbf{Phạm vi}} \\
\midrule
\rowcolor{rowgray} 1 & JPEG Quality & 10--63 (thấp = chất lượng cao hơn) \\
2 & Frame Size & QQVGA đến UXGA (14 mức) \\
\rowcolor{rowgray} 3--5 & Brightness, Contrast, Saturation & $-2$ đến $+2$ \\
6--9 & AWB, AEC, AGC, AEC2 & 0/1 (bật/tắt) \\
\rowcolor{rowgray} 16--17 & H-Mirror, V-Flip & 0/1 \\
20 & XCLK Frequency & 6/8/12/16/24 MHz \\
\rowcolor{rowgray} 21 & Frame Buffer Count & 1--4 \\
\bottomrule
\end{tabular}
\end{table}

ESP32 phản hồi cài đặt hiện tại qua packet với magic header \texttt{0xBEEF}
theo format: \texttt{[0xBE][0xEF][ID][Value]...} cho tất cả 21 tham số.
Phản hồi này được gửi mỗi khi có thay đổi cài đặt, giúp giao diện web
luôn đồng bộ với trạng thái thực tế của camera.

\implDenoise
\implVoice
\implClassify
